\chapter{Interaction with multiflow}
\label{ch:multiflow}

\multi{} uses \mftl{} for its front-tracking interfaces and immersed-boundary bodies. This
chapter describes the interface between the two as it stands after Stage~16; the details live in
\multi{} (\texttt{src/general/mftlwrapper/}, \texttt{src/include/mftl\_adapter*.h} and
\texttt{src/interfacial/fronttracking/}).

\section{The adapter}

\multi{} never calls \mftl{} directly from its physics code. An adapter layer
(\texttt{mftl\_adapter.h}, implemented in \texttt{mftl\_adapter\_mftl2.c}) offers the operations
the solver needs --- body lists, iteration over triangles and vertices, positions, geometry in
local frames, redistribution, remeshing, I/O --- and maps them onto \mftl:
\begin{itemize}
  \item the decomposition of \multi's front-tracking overlap mesh (its particle-box grid)
    becomes the domain of an ensemble created with \fn{MFTLEnsembleCreateRectilinear}, with the
    per-axis cuts and the rank map of that grid. It is aligned with, but not identical to, the
    fluid partition; \multi's \texttt{sendnonlocal*} exchanges reconcile the two;
  \item per-body operations that \mftl{} implements for ensembles (remeshing, serial brackets)
    run on a one-object view of the body's ensemble;
  \item every redistribution first refreshes the cached centres of the owned triangles (the
    ownership rule of \S\ref{sec:triprops});
  \item physics code reads positions only through adapter accessors that return local-frame or
    continuous coordinates (\texttt{MFAdapterTriVertexPositions}, \texttt{MFAdapterVertContinuous},
    \texttt{MFAdapterVertDelta}, \texttt{MFAdapterEdgeVector}), so no physics routine subtracts
    wrapped positions itself; after advection \texttt{MFAdapterBodyWrapVertices} applies the
    displacement rule;
  \item checked queries (volume, centre, inside/outside) return
    \texttt{MFADAPTER\_DEFINED}, \texttt{\_UNDEFINED} or \texttt{\_UNSUPPORTED}, mirroring the
    phase-set statuses.
\end{itemize}
\multi{} finds \mftl{} with \texttt{find\_package(MFTL2)} (its CMake variable \texttt{MFTL2\_ROOT})
and needs a LAPACK-enabled build for its default vertex fit.

\section{Declaring periodic and open fronts}

The front-tracking options that concern \mftl{} (all in the \texttt{FRONTTRACKING} section of the
options file) are explicit declarations; nothing is inferred from the geometry:
\begin{center}\small
\begin{tabularx}{\textwidth}{@{}lX@{}}
\toprule
Option & Meaning \\
\midrule
\texttt{PeriodicFronts XZ} & the fronts are periodic along the listed axes (each must be periodic in the fluid) \\
\texttt{BodyMesh <body> <file.gts>} & read the body from a GTS file instead of generating a sphere \\
\texttt{OpenSurface <body> BoundaryFromMesh} & the body is open; every unmatched edge of its mesh is a boundary, tagged per loop \\
\texttt{RemeshMode <body> GENERALIZED|DEFAULT} & opt a closed body into the generalized remesher \\
\texttt{SurfaceScalarBC <body> <tag|ALL> NEUMANN q} & scalar flux $q$ per unit length (positive outward along $\vect{\nu}$) at the loops with that tag \\
\texttt{SurfaceScalarBC <body> <tag|ALL> DIRICHLET c} & prescribed concentration at the loops \\
\texttt{SurfaceTensionBoundary <body> <tag|ALL> PINNED STATIC} & the loops are pinned to a static support \\
\texttt{SurfaceTensionBoundary <body> <tag|ALL> REJECT} & (default) refuse surface tension on this open body \\
\texttt{PhaseSetMember <body> <+1|-1>} & the open body is a member of the dispersed-phase set with this sign \\
\texttt{WallClosure <X|Y|Z> <MIN|MAX>} & the dispersed-phase set is closed at this (non-periodic) fluid wall \\
\bottomrule
\end{tabularx}
\end{center}
\texttt{ROBIN} (scalar) and \texttt{CONTACT\_LINE} and IBM/DEM supports (surface tension) are
reserved and rejected.

\section{The capability check}

At initialisation --- and again after every class transition --- \multi{} classifies every front
body as \emph{closed} (non-periodic fronts), \emph{closed periodic compact}, \emph{closed periodic
wrapping} or \emph{open}, and checks it against every enabled module. A conflict stops the run with
a message that names the module and the body. The rules:
\begin{itemize}
  \item surface-only modules (Frenet-element surface tension, surface scalar transport) run on
    every closed class;
  \item modules that need an enclosed phase (indicator function, CSF, integral surface tension,
    marching cubes with breakup and coalescence, volume recovery, force computation, generic
    acceleration, normal-corrected advection, IRL reconstruction) need a closed non-periodic
    front, or a periodic front whose dispersed-phase set is \texttt{MFTL\_DEFINED};
  \item on periodic fronts some modules are not yet implemented (CSF, integral surface tension,
    surface viscosity, normal-corrected advection, generic acceleration, force computation,
    divergence-free interpolation, IRL reconstruction); the phase modules also reject wrapping
    fronts;
  \item an open body runs a module only with the declarations the module needs: scalar
    transport needs a scalar boundary condition for every boundary tag (zero flux is never
    assumed); Frenet surface tension needs \texttt{PINNED STATIC} on every tag; the indicator
    and volume recovery need membership in a wall-closed, \texttt{DEFINED} dispersed-phase set whose
    loops are all pinned; CSF and IRL additionally need a non-periodic front and a two-ring
    fit stencil; integral surface tension, marching cubes and surface viscosity are rejected.
\end{itemize}

\section{Boundary conditions on open bodies}
\label{sec:mf-bcs}

\mftl{} supplies only geometry for boundary edges: tag, length $L_e$, tangent $\vect{t}$, conormal
$\vect{\nu}$, midpoint and the triangle height $h_T=2A_T/L_e$
(\fn{MFTLBoundaryEdgeGeometry}, \fn{MFTLBoundaryEdgeCentroidDistance}). \multi{} builds its
boundary conditions from them.

\paragraph{Surface scalar.} The scalar lives at triangle centroids; the distance from the centroid
to a boundary edge along $\vect{\nu}$ is $\delta_e=h_T/3$.
\begin{itemize}
  \item \texttt{NEUMANN q}: the face value in the gradient reconstruction is
    $c_f=c_T-q\,\delta_e/D$, the diffusion right-hand side gets $-qL_e$, and the mass changes by
    $-\sum qL_e\Delta t$ (zero for $q=0$);
  \item \texttt{DIRICHLET c}: the face value is $c_D$, and the two-point flux
    $D L_e(c_D-c_T)/\delta_{\text{eff}}$ uses $\delta_{\text{eff}}=\max(\delta_e,0.1L_e)$ (a floor
    that bounds the coefficient on slivers; every use is counted).
\end{itemize}
The advective flux across a boundary edge is zero (a material boundary).

\paragraph{Pinned surface tension.} For constant $\sigma$ the element force of a triangle is
$\sum_{\text{edges}}\sigma L\vect{\nu}_e$; summed over an open surface the interior edges cancel,
leaving the line force at the boundary. With \texttt{PINNED STATIC} the boundary vertices are held
on their support (their loop is a frozen, constrained boundary curve), the element forces are
spread to the fluid unchanged, and the line force
\begin{equation}
  \vect{R} = -\sum_{\text{loop}} \sigma L_e\,\vect{\nu}_e
\end{equation}
is the reaction on the support, reported per step: the fluid force plus $\vect{R}$ is zero to
round-off. A free boundary line force is never applied implicitly.

\section{Periodic fronts}

With \texttt{PeriodicFronts} the \mftl{} domain is periodic on the declared axes with the fluid
period. Peskin spreading and interpolation use wrapped cell ranges, VTU output draws seam
triangles per image, logs report continuous centroids on compact axes (and ``undefined''
otherwise), and restarts use native mesh version~2. Closed bodies keep the \mftlone-parity remesher
whether periodic or not; open bodies always use the generalized remesher.
